% Part: incompleteness
% Chapter: arithmetization-syntax
% Section: proofs-in-lk

\documentclass[../../../include/open-logic-section]{subfiles}

\begin{document}

\olfileid{inc}{art}{plk}
\olsection{$\Log{LK}$లో \usetoken{P}{derivation}}

\begin{explain}
\tetoken{వ్యుత్పత్తులను}{!!{derivation}s} అంకగణితీకరించాలంటే,
\tetoken{వ్యుత్పత్తులను}{!!{derivation}s} సంఖ్యలుగా ప్రాతినిధ్యం చేయాలి.
\tetoken{వ్యుత్పత్తులు}{!!{derivation}s}
సీక్వెంట్ల వృక్షాలు, ప్రతి నిగమనానికీ ఒక గుర్తింపు చీటీ కూడా ఉంటుంది కనుక,
పునరావృత్త ప్రాతినిధ్యమే అత్యంత సహజమైన పద్ధతి. ఒక
\tetoken{వ్యుత్పత్తిని}{!!{derivation}} అంత్య-సీక్వెంట్, గుర్తింపు చీటీ,
చివరి నిగమనపు పూర్వాధారాలకు దారితీసే ఉప-\tetoken{వ్యుత్పత్తుల}{!!{derivation}s}
ప్రాతినిధ్యాలు అనే భాగాలు గల క్రమిత సమూహంగా ప్రాతినిధ్యం చేస్తాం.
\end{explain}

\begin{defn}
$\Gamma$ అనేది \tetoken{వాక్యాల}{!!{sentence}s} పరిమిత క్రమం,
$\Gamma=\tuple{!A_1,\dots,!A_n}$ అయితే,
$\Gn{\Gamma}=\tuple{\Gn{!A_1},\dots,\Gn{!A_n}}$.

$\Gamma\Sequent\Delta$ ఒక సీక్వెంట్ అయితే, $\Gamma\Sequent\Delta$ యొక్క
గ్యోడెల్ సంఖ్య
\[
\Gn{\Gamma \Sequent \Delta}=\tuple{\Gn{\Gamma},\Gn{\Delta}}
\]

$\pi$ అనేది $\Log{LK}$లో \tetoken{ఒక వ్యుత్పత్తి}{!!a{derivation}} అయితే,
$\Gn{\pi}$ను ఇలా నిర్వచిస్తాం:
\begin{enumerate}
\item $\pi$లో ప్రారంభ సీక్వెంట్ $\Gamma\Sequent\Delta$ మాత్రమే ఉంటే,
  $\Gn{\pi}$ అనేది
  \[
    \tuple{0,\Gn{\Gamma\Sequent\Delta}}.
  \]
\item $\pi$ ఒకటి లేదా రెండు పూర్వాధారాలు గల నిగమనంతో ముగిసి,
  $\Gamma\Sequent\Delta$ దాని నిష్కర్షగా ఉంటే, $\pi_1$, $\pi_2$లు చివరి
  నిగమనపు పూర్వాధారాల వద్ద ముగిసే తక్షణ ఉపనిరూపణలు అయితే, వరుసగా
  $\Gn{\pi}$ అనేది
  \begin{align*}
    & \tuple{1,\Gn{\pi_1},\Gn{\Gamma\Sequent\Delta},k}\text{ లేదా}\\
    & \tuple{2,\Gn{\pi_1},\Gn{\pi_2},\Gn{\Gamma\Sequent\Delta},k},
  \end{align*}
  చివరి నిగమనంలో వాడిన నియమాన్ని బట్టి $k$ను ఈ పట్టిక ఇస్తుంది:

  \begin{tabular}{lccccccc}
    \text{నియమం:} & \LeftR{\Weakening} & \RightR{\Weakening} &
    \LeftR{\Contraction} & \RightR{\Contraction} &
    \LeftR{\Exchange} & \RightR{\Exchange} \\
    $k$: & 1 & 2 & 3 & 4 & 5 & 6 \\[2ex]
    \text{నియమం:} & \LeftR{\lnot} & \RightR{\lnot} &
    \LeftR{\land} & \RightR{\land} &
    \LeftR{\lor} & \RightR{\lor} \\
    $k$: & 7 & 8 & 9 & 10 & 11 & 12 \\[2ex]
    \text{నియమం:} & \LeftR{\lif} & \RightR{\lif} &
    \LeftR{\lforall} & \RightR{\lforall} &
    \LeftR{\lexists} & \RightR{\lexists} \\
    $k$: & 13 & 14 & 15 & 16 & 17 & 18 \\[2ex]
    \text{నియమం:} & \Cut & = \\
    $k$: & 19 & 20
  \end{tabular}
\end{enumerate}
\end{defn}

\begin{ex}
  ఈ అతి సరళమైన \tetoken{వ్యుత్పత్తిని}{!!{derivation}} పరిగణించండి:
  \begin{prooftree}
    \Axiom$!A \fCenter !A$
    \RightLabel{\LeftR{\land}}
    \UnaryInf$!A \land !B \fCenter !A$
    \RightLabel{\RightR{\lif}}
    \UnaryInf$\fCenter (!A \land !B) \lif !A$
  \end{prooftree}
  ప్రారంభ సీక్వెంట్ గ్యోడెల్ సంఖ్య
  $p_0=\tuple{0,\Gn{!A\Sequent !A}}$ అవుతుంది. $\LeftR{\land}$ నిష్కర్షతో
  ముగిసే \tetoken{వ్యుత్పత్తి}{!!{derivation}} గ్యోడెల్ సంఖ్య
  $p_1=\tuple{1,p_0,\Gn{!A\land !B\Sequent !A},9}$ అవుతుంది
  ($\LeftR{\land}$కు ఒక పూర్వాధారం ఉంది కనుక $1$; నిష్కర్ష
  $!A\land !B\Sequent !A$ యొక్క గ్యోడెల్ సంఖ్య; $\LeftR{\land}$ను
  సంకేతీకరించే సంఖ్య $9$). అప్పుడు మొత్తం
  \tetoken{వ్యుత్పత్తి}{!!{derivation}} గ్యోడెల్ సంఖ్య
  $\tuple{1,p_1,\Gn{\Sequent(!A\land !B)\lif !A},14}$, అంటే
  \[
  \tuple{1,\tuple{1,\tuple{0,\Gn{!A\Sequent !A}},\Gn{!A\land !B\Sequent !A},9},
    \Gn{\Sequent(!A\land !B)\lif !A},14}.
  \]
  \sourcecorrection{OLTEART-004}{ఆంగ్ల మూలంలోని సంక్షిప్త గ్యోడెల్ సంకేతంలో $\Gn{\Sequent(!A\land !B)\lif !A)}$ వద్ద, విస్తరించిన క్రమిత సమూహంలో $\Gn{!A\Sequent !A)}$ వద్ద అదనపు కుడి కుండలీకరణ గుర్తులు ఉన్నాయి. పక్కనే ఇచ్చిన నిరూపణ నిష్కర్షకు, ముందే నిర్వచించిన $p_0$కు అనుగుణంగా వాటిని తొలగించాం.}
\end{ex}

\begin{explain}
\tetoken{వ్యుత్పత్తుల}{!!{derivation}s} ప్రాతినిధ్యాన్ని నిర్ణయించిన తరువాత,
అటువంటి \tetoken{వ్యుత్పత్తులను}{!!{derivation}s} ఆదిమ పునరావృత్తంగా
మార్చగలమనీ, వాటి ముఖ్య ధర్మాలు, సంబంధాలను అలాగే వ్యక్తీకరించగలమనీ చూపాలి.
కొన్ని క్రియలు సరళమైనవి: ఉదాహరణకు,
\tetoken{ఒక వ్యుత్పత్తి}{!!a{derivation}} గ్యోడెల్ సంఖ్య~$p$ ఇస్తే,
$\fn{EndSequent}(p)=(p)_{(p)_0+1}$ దాని అంత్య-సీక్వెంట్ గ్యోడెల్ సంఖ్యను,
$\fn{LastRule}(p)=(p)_{(p)_0+2}$ దాని చివరి నియమపు సంకేతసంఖ్యను ఇస్తాయి.
\sourcecorrection{OLTEART-005}{ఆంగ్ల మూలం ఇక్కడ ప్రమేయాన్ని $\fn{EndSeq}$గా నిర్వచించి, మిగతా విభాగమంతా $\fn{EndSequent}$గా ఉపయోగిస్తుంది. ఒకే ప్రమేయం కోసం తరువాత స్థిరంగా వాడిన $\fn{EndSequent}$ పేరును ఇక్కడ కూడా ఇచ్చాం.}
$\fn{Sequent}(s)$ ధర్మాన్ని
\[
  \len{s}=2\land\bforall{i<\len{(s)_0}+\len{(s)_1}}{\fn{Sent}(((s)_0\concat(s)_1)_i)}
\]
అని నిర్వచిస్తే, \tetoken{వాక్యాలతో}{!!{sentence}s} కూడిన సీక్వెంట్‌కు
$s$కు అది నిలవడానికి అవసరమైనదీ, సరిపడేదీ $s$ గ్యోడెల్ సంఖ్య కావడం.
కొన్ని క్రియలు మరింత కష్టం.
వాటిని ఎలా చేయాలో కనీసం రూపురేఖలుగా చూద్దాం. ``$\pi$ అనేది $\Gamma$ నుంచి
$!A$ యొక్క \tetoken{ఒక వ్యుత్పత్తి}{!!a{derivation}}'' అనే సంబంధం
$\pi$, $!A$ల గ్యోడెల్ సంఖ్యల ఆదిమ పునరావృత్త సంబంధమని చూపడమే లక్ష్యం.
\end{explain}

\begin{prop}\ollabel{prop:followsby}
  గ్యోడెల్ సంఖ్య~$p$ గల \tetoken{వ్యుత్పత్తి}{!!{derivation}}~$\pi$లోని
  చివరి నిగమనం సరైనదైతేనే నిలిచే $\fn{Correct}(p)$ ధర్మం ఆదిమ పునరావృత్తం.
\end{prop}

\begin{proof}
  $\Gamma\Sequent\Delta$ ప్రారంభ సీక్వెంట్ కావడానికి అవసరమైనదీ,
  సరిపడేదీ: ఏదో \tetoken{ఒక వాక్యం}{!!a{sentence}}~$!A$ ఉండి
  $\Gamma\Sequent\Delta$ అనేది $!A\Sequent !A$ కావడం, లేదా ఏదో పదం~$t$
  ఉండి $\Gamma\Sequent\Delta$ అనేది $\emptyset\Sequent\eq[t][t]$ కావడం.
  గ్యోడెల్ సంఖ్యల పరంగా,
  క్రింది షరతు నెరవేరినప్పుడే $\fn{InitSeq}(s)$ నిలుస్తుంది:
  \begin{align*}
  \bexists{x<s}{}(\fn{Sent}(x)&\land
  s=\tuple{\tuple{x},\tuple{x}})
  \lor{}\\
  \bexists{t<s}{}(\fn{Term}(t)&\land
  s=\tuple{0,\tuple{\Gn{{\eq}(}\concat t\concat\Gn{,}\concat t\concat\Gn{)}}}).
  \end{align*}

  ప్రతి నిగమన నియమం~$R$కి $\fn{FollowsBy}_R(p)$ సంబంధం ఆదిమ పునరావృత్తమని
  కూడా చూపాలి. $p$ అనేది \tetoken{వ్యుత్పత్తి}{!!{derivation}}~$\pi$కి
  గ్యోడెల్ సంఖ్యగా ఉండి, $\pi$ యొక్క అంత్య-సీక్వెంట్ దాని తక్షణ
  ఉప-\tetoken{వ్యుత్పత్తుల}{!!{derivation}s} నుంచి $R$ను సరిగా
  వర్తింపజేయడం వల్ల వచ్చినప్పుడే $\fn{FollowsBy}_R(p)$ నిలుస్తుంది.

  సరళమైన సందర్భం \RightR{\land} నియమం. $\pi$ సరైన
  $\RightR{\land}$ నిగమనంతో ముగిస్తే అది ఇలా ఉంటుంది:
  \begin{prooftree}
    \AxiomC{}
    \RightLabel{$\pi_1$}
    \Deduce$\Gamma \fCenter \Delta, !A$

    \AxiomC{}
    \RightLabel{$\pi_2$}
    \Deduce$\Gamma \fCenter \Delta, !B$

    \RightLabel{\RightR\land}
    \BinaryInf$\Gamma \fCenter \Delta, !A \land !B$
  \end{prooftree}
  కాబట్టి, $\Gamma$, $\Delta$ అనే \tetoken{వాక్యాల}{!!{sentence}s} క్రమాలు,
  $!A$, $!B$ అనే రెండు \tetoken{వాక్యాలు}{!!{sentence}s} ఉండి,
  $\pi_1$ అంత్య-సీక్వెంట్ $\Gamma\Sequent\Delta,!A$,
  $\pi_2$ అంత్య-సీక్వెంట్ $\Gamma\Sequent\Delta,!B$, $\pi$ అంత్య-సీక్వెంట్
  $\Gamma\Sequent\Delta,!A\land !B$ అయితేనే
  \tetoken{వ్యుత్పత్తి}{!!{derivation}}~$\pi$లోని చివరి నిగమనం
  $\RightR{\land}$ యొక్క సరైన వర్తింపు. దీన్ని గ్యోడెల్ సంఖ్యలలోకి
  అనువదిస్తే చాలు. $s=\Gn{\Gamma\Sequent\Delta}$ అయితే
  $(s)_0=\Gn{\Gamma}$, $(s)_1=\Gn{\Delta}$. అందువల్ల
  $\fn{FollowsBy}_{\RightR{\land}}(p)$ నిలవడానికి అవసరమైనదీ, సరిపడేదీ
\begin{align*}
& \bexists{g<p}{\bexists{d<p}{\bexists{a<p}{\bexists{b<p}{\quad}}}}\\
& \qquad \fn{EndSequent}(p)=
  \tuple{g,d\concat
    \tuple{\Gn{(}\concat a\concat\Gn{\land}\concat b\concat\Gn{)}}}\land{}\\
& \qquad \fn{EndSequent}((p)_1)=\tuple{g,d\concat\tuple{a}}\land{}\\
& \qquad \fn{EndSequent}((p)_2)=\tuple{g,d\concat\tuple{b}}\land{}\\
& \qquad (p)_0=2\land\fn{LastRule}(p)=10.
\end{align*}
వరుస పంక్తులు ఇలా చెబుతాయి: ``గ్యోడెల్ సంఖ్య~$g$ గల క్రమం~($\Gamma$),
గ్యోడెల్ సంఖ్య~$d$ గల క్రమం~($\Delta$), గ్యోడెల్ సంఖ్య~$a$ గల
\tetoken{ఒక సూత్రం}{!!a{formula}}~($!A$), గ్యోడెల్ సంఖ్య~$b$ గల
\tetoken{ఒక సూత్రం}{!!a{formula}}~($!B$) ఉన్నాయి''; ``$\pi$ అంత్య-సీక్వెంట్
$\Gamma\Sequent\Delta,!A\land !B$''; ``$\pi_1$ అంత్య-సీక్వెంట్
$\Gamma\Sequent\Delta,!A$''; ``$\pi_2$ అంత్య-సీక్వెంట్
$\Gamma\Sequent\Delta,!B$''; ``$\pi$కి రెండు తక్షణ ఉపవ్యుత్పత్తులు ఉన్నాయి,
చివరి నిగమన నియమం $\RightR\land$ (దాని సంఖ్య~$10$).''

$\pi$లోని చివరి నిగమనం $\RightR\lexists$ యొక్క సరైన వర్తింపు కావడానికి
అవసరమైనదీ, సరిపడేదీ: $\Gamma$, $\Delta$ క్రమాలు,
\tetoken{ఒక సూత్రం}{!!a{formula}}~$!A$, చరం~$x$, పదం~$t$ ఉండి,
$\pi$ అంత్య-సీక్వెంట్ $\Gamma\Sequent\Delta,\lexists[x][!A]$ కావడం,
$\pi_1$ అంత్య-సీక్వెంట్ $\Gamma\Sequent\Delta,\Subst{!A}{t}{x}$ కావడం.
కాబట్టి గ్యోడెల్ సంఖ్యల పరంగా $\fn{FollowsBy}_{\RightR\lexists}(p)$
నిలవడానికి అవసరమైనదీ, సరిపడేదీ
\begin{align*}
  & \bexists{g<p}{\bexists{d<p}{\bexists{a<p}{\bexists{x<p}{\bexists{t<p}{\quad}}}}}\\
  & \qquad \fn{EndSequent}(p)=
  \tuple{
    g,
    d\concat\tuple{\Gn{\lexists}\concat x\concat a}
  }\land{}\\
  & \qquad \fn{EndSequent}((p)_1)=
  \tuple{
    g,
    d\concat\tuple{\fn{Subst}(a,t,x)}
  }\land{}\\
  & \qquad (p)_0=1\land\fn{LastRule}(p)=18.
\end{align*}
తరువాత $\fn{Correct}(p)$ను ఇలా నిర్వచిస్తాం:
\begin{multline*}
  \fn{Sequent}(\fn{EndSequent}(p))\land{}\\
  [(\fn{LastRule}(p)=1\land
  \fn{FollowsBy}_{\LeftR\Weakening}(p))\lor\dots\lor{}\\
  (\fn{LastRule}(p)=20\land\fn{FollowsBy}_{\eq}(p))\lor{}\\
  (p)_0=0\land\fn{InitSeq}(\fn{EndSequent}(p))]
\end{multline*}
\sourcecorrection{OLTEART-006}{ఆంగ్ల మూలం ప్రారంభ సీక్వెంట్ పరీక్షను $\fn{InitSeq}$గా నిర్వచించి, చివరి సూత్రంలో $\fn{InitialSeq}$గా పిలుస్తుంది. నిర్వచించిన $\fn{InitSeq}$ పేరును స్థిరంగా వాడాం.}
మొదటి పంక్తి $p$ యొక్క అంత్య-సీక్వెంట్ నిజంగా
\tetoken{వాక్యాలతో}{!!{sentence}s} కూడిన సీక్వెంట్ అని నిర్ధారిస్తుంది.
చివరి పంక్తి $p$ కేవలం ప్రారంభ సీక్వెంట్ అయిన సందర్భాన్ని కలుపుతుంది.
\sourcecorrection{OLTEART-007}{ఆంగ్ల మూలంలోని ఈ వివరణ మొదటి పంక్తి ``$d$ యొక్క అంత్య-సీక్వెంట్''ను పరీక్షిస్తుందని అంటుంది; సూత్రం, ఈ విభాగపు చరం రెండూ $p$ను వాడుతున్నందున దాన్ని $p$గా సరిచేశాం.}
\end{proof}

\begin{prob}
  \olref[inc][art][plk]{prop:followsby}లోలాగే ఈ ధర్మాలను నిర్వచించండి:
  \begin{enumerate}
  \item $\fn{FollowsBy}_{\Cut}(p)$,
  \item $\fn{FollowsBy}_{\LeftR\lif}(p)$,
  \item $\fn{FollowsBy}_{\eq}(p)$,
  \item $\fn{FollowsBy}_{\RightR{\lforall}}(p)$.
  \end{enumerate}
  చివరిదానికి, గ్యోడెల్ సంఖ్య~$p$ గల
  \tetoken{వ్యుత్పత్తి}{!!{derivation}} యొక్క చివరి నిగమనం ఐగెన్ చరరాశి
  షరతును నెరవేరుస్తుందా, అంటే $\RightR{\lforall}$ యొక్క ఐగెన్
  చరరాశి~$a$ అంత్య-సీక్వెంట్‌లో కనిపించదా అని ఆదిమ పునరావృత్తంగా
  పరీక్షించగలమని కూడా చూపాలి.
\end{prob}

\begin{prop}
  \ollabel{prop:deriv}
  $p$ సరైన \tetoken{వ్యుత్పత్తి}{!!{derivation}}~$\pi$ యొక్క గ్యోడెల్ సంఖ్య
  అయితే నిలిచే $\fn{Deriv}(p)$ సంబంధం ఆదిమ పునరావృత్తం.
\end{prop}

\begin{proof}
  \tetoken{ఒక వ్యుత్పత్తి}{!!^a{derivation}}~$\pi$లోని ప్రతి నిగమనం ఒక
  నియమపు సరైన వర్తింపు అయితే, అంటే దాని ప్రతి ఉప-\tetoken{వ్యుత్పత్తి}{!!{derivation}s}
  సరైన నిగమనంతో ముగిస్తే, $\pi$ సరైనది. కాబట్టి $\fn{Deriv}(p)$ నిలవడానికి
  అవసరమైనదీ, సరిపడేదీ
  \[
  \bforall{i<\len{\fn{SubtreeSeq}(p)}}{\fn{Correct}((\fn{SubtreeSeq}(p))_i)}.
  \]
  \sourcecorrection{OLTEART-014}{ఆంగ్ల మూలంలోని ప్రదర్శిత సూత్రంలో $\fn{Correct}$ ఆర్గుమెంటును ముగించే కుడి కుండలీకరణ గుర్తు లేదు. విధేయ అన్వయాన్ని సరిగ్గా ముగించే గుర్తును చేర్చాం.}
  \sourcecorrection{OLTEART-008}{ఆంగ్ల మూలంలోని ప్రతిపాదన $\fn{Deriv}(p)$ను పేర్కొన్నప్పటికీ, నిరూపణ వాక్యం ఒక్కసారిగా $\fn{Deriv}(d)$కు మారి, వెంటనే వచ్చే సూత్రం మళ్లీ $p$ను వాడుతుంది. ప్రతిపాదన, సూత్రానికి అనుగుణంగా నిరూపణ వాక్యంలోనూ $p$ను వాడాం.}
\end{proof}


\begin{prop}
$\Gamma$ అనేది \tetoken{వాక్యాల}{!!{sentence}s} ఆదిమ పునరావృత్త సమితి
అనుకుందాం. ``ఏదో పరిమిత $\Gamma_0\subseteq\Gamma$కు
$\Gamma_0\Sequent !A$ యొక్క \tetoken{ఒక వ్యుత్పత్తి}{!!a{derivation}}~$\pi$
సంకేతసంఖ్య $x$, $!A$ గ్యోడెల్ సంఖ్య $y$'' అని వ్యక్తీకరించే
$\Prf[\Gamma](x,y)$ సంబంధం ఆదిమ పునరావృత్తం.
\end{prop}

\begin{proof}
``$y\in\Gamma$''ను ఆదిమ పునరావృత్త విధేయం $R_\Gamma(y)$ ఇస్తుందని
అనుకుందాం. $y$ ఒక వాక్యం~$!A$ గ్యోడెల్ సంఖ్యగా, $x$ అంత్య-సీక్వెంట్
$\Gamma_0\Sequent !A$ గల $\Log{LK}$-\tetoken{వ్యుత్పత్తి}{!!{derivation}}
సంకేతసంఖ్యగా ఉన్నప్పుడే నిలిచే $\Prf[\Gamma](x,y)$ ఆదిమ పునరావృత్తమని
చూపాలి.

మునుపటి ప్రతిపాదన ప్రకారం, $x$ అనేది $\Log{LK}$లో సరైన
\tetoken{వ్యుత్పత్తి}{!!{derivation}}~$\pi$ సంకేతసంఖ్య అయినప్పుడే నిలిచే
$\fn{Deriv}(x)$ ధర్మం ఆదిమ పునరావృత్తం. $x$ అటువంటి సంకేతసంఖ్య అయితే,
$\fn{EndSequent}(x)$ అనేది $\pi$ అంత్య-సీక్వెంట్ సంకేతసంఖ్య;
కాబట్టి $(\fn{EndSequent}(x))_0$ అంత్య-సీక్వెంట్ ఎడమ వైపు సంకేతసంఖ్య,
$(\fn{EndSequent}(x))_1$ కుడి వైపు సంకేతసంఖ్య. అందువల్ల ``$\pi$
అంత్య-సీక్వెంట్ కుడి వైపు $!A$'' అని
$\len{(\fn{EndSequent}(x))_1}=1\land((\fn{EndSequent}(x))_1)_0=y$
గా వ్యక్తీకరించవచ్చు.
\sourcecorrection{OLTEART-009}{ఆంగ్ల మూలం ఈ మధ్యంతర సమీకరణలో కుడి వైపు ఏకైక సూత్రపు గ్యోడెల్ సంఖ్యను $x$తో సమానం చేస్తుంది; $x$ వ్యుత్పత్తి సంకేతసంఖ్య, $y$నే $!A$ గ్యోడెల్ సంఖ్య. వెంటనే వచ్చే తుది నిర్వచనానికి అనుగుణంగా $y$గా సరిచేశాం.}
$\pi$ అంత్య-సీక్వెంట్ ఎడమ వైపు సహజంగానే పరిమితం; దానిలోని ప్రతి వాక్యం
$\Gamma$లో ఉందని మాత్రమే వ్యక్తీకరించాలి. కాబట్టి $\Prf[\Gamma](x,y)$ను
ఇలా నిర్వచించవచ్చు:
\begin{align*}
\Prf[\Gamma](x,y)\defiff{}&
\fn{Deriv}(x)\land{}\\
&\bforall{i<
  \len{(\fn{EndSequent}(x))_0}}{R_\Gamma(((\fn{EndSequent}(x))_0)_i)}\land{}\\
&\len{(\fn{EndSequent}(x))_1}=1\land((\fn{EndSequent}(x))_1)_0=y.
\end{align*}
\end{proof}

\end{document}
